\chapter{Blocks 9--10: Applications --- Single-Cell \& GWAS}\label{chap:apps}

\textbf{Primary:} scRNA-seq pipeline, GWAS pipeline\quad|\quad\textbf{Interleave:} All previous skills converge here

\begin{brainbox}[Where it all comes together]
Unix, Python, R, and statistics all appear in these two blocks. You will run complete analyses from raw data to biological interpretation.
\end{brainbox}


\section{Block 9: Single-Cell RNA-seq, Start to Finish}

\sprint{1}{Technology and Data Orientation}{25}
\begin{itemize}[nosep]
  \item 10x Chromium captures RNA from individual cells, producing a UMI count matrix with cells as rows and genes as columns in \texttt{AnnData}.
  \item Download a public PBMC dataset from the 10x Genomics website (\url{https://www.10xgenomics.com/datasets}); a pre-processed \texttt{.h5ad} is fine to start.
  \item Install scanpy if needed: \texttt{mamba install -c conda-forge -c bioconda scanpy}. Docs: \url{https://scanpy.readthedocs.io/}. Load with \texttt{adata = sc.read\_h5ad('pbmc.h5ad')}.
  \item Explore: \texttt{adata.shape}, \texttt{adata.obs.head()}, \texttt{adata.var.head()}.
  \item Answer in your journal: what does each row of \texttt{adata} represent? Each column? What is a plausible range for per-cell total counts and for number of expressed genes? (For a healthy PBMC droplet library, 500--10{,}000 genes and 1{,}000--50{,}000 UMIs per cell are common.)
\end{itemize}
\begin{winbox}[Checkpoint]
You loaded a real single-cell dataset and can describe its dimensions and what each axis means.
\end{winbox}

\sprint{2}{Quality Control: Filtering Cells and Doublets}{25}
\begin{itemize}[nosep]
  \item Flag mitochondrial genes so they can be counted:
    \begin{lstlisting}
adata.var['mt'] = adata.var_names.str.startswith('MT-')
    \end{lstlisting}
  \item Calculate QC metrics:
    \begin{lstlisting}
sc.pp.calculate_qc_metrics(adata, qc_vars=['mt'])
    \end{lstlisting}
    This adds \texttt{n\_genes\_by\_counts}, \texttt{total\_counts}, and \texttt{pct\_counts\_mt} to \texttt{adata.obs}.
  \item Plot distributions:
    \begin{lstlisting}
sc.pl.violin(adata,
  ['n_genes_by_counts',
   'total_counts',
   'pct_counts_mt'])
    \end{lstlisting}
  \item Filter: keep cells with roughly 200--5000 detected genes and fewer than 20\% mitochondrial reads. The biological rationale: very few genes suggests empty droplets or dying cells; very many suggests doublets; high mitochondrial fraction suggests cells whose membranes have leaked cytoplasmic RNA.
  \item Doublet detection with Scrublet: \texttt{sc.pp.scrublet(adata)} (or \texttt{sc.external.pp.scrublet(adata)} on older scanpy versions).
\end{itemize}
\begin{winbox}[Checkpoint]
You filtered out low-quality cells and doublets. Record cell counts before and after in \texttt{journal.md}. A typical PBMC dataset loses 10--30\% of droplets.
\end{winbox}

\sprint{3}{Normalize $\to$ HVGs $\to$ PCA $\to$ UMAP $\to$ Cluster}{25}
\begin{itemize}[nosep]
  \item \texttt{sc.pp.normalize\_total(adata); sc.pp.log1p(adata)} --- library size normalization
  \item \texttt{sc.pp.highly\_variable\_genes(adata)} --- select $\sim$2000 informative genes
  \item Run the pipeline:
    \begin{lstlisting}
sc.tl.pca(adata)
sc.tl.neighbors(adata)
sc.tl.umap(adata)
sc.tl.leiden(adata, resolution=0.5)
    \end{lstlisting}
  \item \texttt{sc.pl.umap(adata, color=['leiden'])} --- a single-cell cluster plot from raw data
\end{itemize}
\begin{winbox}[Checkpoint]
You went from filtered counts to a clustered UMAP plot with six scanpy calls. Save the plot to your repo.
\end{winbox}

\sprint{4}{Annotate Cell Types and Pseudobulk DE}{25}
\begin{itemize}[nosep]
  \item Plot canonical markers: \texttt{sc.pl.dotplot(adata, ['CD3D', 'CD14', 'MS4A1', 'NKG7'], groupby='leiden')}. For PBMCs: \texttt{CD3D} marks T cells, \texttt{CD14} monocytes, \texttt{MS4A1} B cells, \texttt{NKG7} NK/cytotoxic cells.
  \item Assign labels based on which clusters light up for which markers:
    \begin{lstlisting}
adata.obs['cell_type'] = adata.obs['leiden'].map(
  {'0': 'CD4 T', '1': 'Monocytes', ...})
    \end{lstlisting}
  \item For differential expression across conditions, use pseudobulk: aggregate UMI counts per (donor, cell\_type) pair and feed the resulting count matrix to DESeq2 or edgeR in R.
  \item Why pseudobulk? Cell-level tests treat every cell as an independent observation, but cells from the same donor share many sources of variation. Pseudobulk respects the actual unit of replication (the donor) and produces calibrated p-values.
\end{itemize}
\begin{winbox}[Checkpoint]
You annotated cell types from marker expression, ran a pseudobulk DE test, and can state in one sentence why pseudobulk gives more honest p-values than per-cell testing.
\end{winbox}

\begin{tipbox}[Multi-modal extensions (for your journal)]
\begin{itemize}[nosep]
  \item CITE-seq: same cells, RNA + protein measurements $\to$ Seurat WNN for joint embedding
  \item scATAC-seq: chromatin accessibility per cell $\to$ ArchR, Signac
  \item Spatial transcriptomics: gene expression + spatial location $\to$ Squidpy, cell2location
  \item Perturb-seq: CRISPR perturbations linked to transcriptomes $\to$ Mixscape, scMAGeCK
\end{itemize}
\end{tipbox}


\section{Block 10: GWAS Pipeline, Variant to Function}

\sprint{5}{Genotype QC with PLINK2}{25}
\begin{itemize}[nosep]
  \item PLINK2 is the workhorse. Install via bioconda (\texttt{mamba install -c bioconda plink2}) or download a pre-built binary from \url{https://www.cog-genomics.org/plink/2.0/}. Get a small dataset, e.g., a HapMap or 1000 Genomes subset in PLINK format.
  \item Sample QC flags: \texttt{--mind} (per-sample call rate), \texttt{--het} (heterozygosity), \texttt{--check-sex}, \texttt{--king-cutoff} (relatedness).
  \item Variant QC flags: \texttt{--geno} (per-variant call rate), \texttt{--hwe} (Hardy--Weinberg equilibrium), \texttt{--maf} (minor allele frequency threshold).
  \item PCA: \texttt{plink2 --pca 10} computes the top 10 principal components. Plot PC1 vs.\ PC2.
\end{itemize}
\begin{winbox}[Checkpoint]
You QC'd a genotype dataset and can report how many samples and variants were dropped by each filter. Every GWAS starts here.
\end{winbox}

\sprint{6}{Association Testing}{25}
\begin{itemize}[nosep]
  \item Simulate a phenotype correlated with a chosen causal variant so you know the true answer and can check your pipeline recovers it.
  \item Run \texttt{plink2 --glm} to fit a linear (quantitative) or logistic (case/control) model at every variant.
  \item For real human data, use REGENIE or BOLT-LMM (mixed models that handle population structure and relatedness simultaneously).
  \item Make a Manhattan plot and a QQ plot of the results, and check the genomic inflation factor $\lambda_\text{GC}$. Values near 1.00 are healthy; values above 1.1--1.2 suggest uncontrolled structure or miscalibration.
\end{itemize}
\begin{winbox}[Checkpoint]
You ran a GWAS on simulated data and recovered the known causal variant as the top hit. Your QQ plot and $\lambda_\text{GC}$ indicate a calibrated test.
\end{winbox}

\sprint{7}{Post-GWAS: Heritability, Fine-Mapping, Interpretation}{25}
\begin{itemize}[nosep]
  \item LDSC (\url{https://github.com/bulik/ldsc}) estimates SNP-heritability from summary statistics: what fraction of phenotypic variance is explained by the common-variant tagging set? Install via conda; follow the repo's instructions for its Python-2 dependency, or use the Python-3 fork \texttt{ldsc-py3}.
  \item Fine-mapping: the SuSiE R package (\url{https://stephenslab.github.io/susieR/}) reports 95\% credible sets, the smallest sets of variants likely to contain the causal one in each signal. \texttt{install.packages("susieR")}.
  \item Colocalization: the \texttt{coloc} R package (\texttt{install.packages("coloc")}, \url{https://chr1swallace.github.io/coloc/}) tests whether your GWAS signal shares a causal variant with a nearby eQTL.
  \item Polygenic risk scores (PRS) aggregate small effects across the genome into a single per-sample predictor, typically by weighted sum of effect alleles.
\end{itemize}
\begin{winbox}[Checkpoint]
For one locus in your GWAS, you reported a heritability estimate for the trait, a fine-mapped credible set, and a colocalization posterior with an eQTL (or a clear ``no colocalization'').
\end{winbox}

\begin{restbox}[End of Blocks 9--10]
You've now run the two most important analysis pipelines in modern genomics end-to-end. Take a break before the final unit.
\end{restbox}

\subsection{Checkpoint:}
\begin{itemize}[nosep]
  \item You can run a complete scRNA-seq analysis from raw counts to annotated clusters
  \item You can explain every step in the single-cell pipeline and why each matters
  \item You can QC genotype data, run an association test, and make Manhattan/QQ plots
  \item You can explain heritability, fine-mapping, and colocalization at a conceptual level
\end{itemize}

\subsection{Resources}
\begin{itemize}[nosep]
  \item Scanpy tutorials (scanpy.readthedocs.io) --- the PBMC 3k tutorial is the canonical starting point
  \item Heumos et al.\ -- Best practices for single-cell analysis (Nat Rev Genet, 2023)
  \item Uffelmann et al.\ -- Genome-wide association studies (Nat Rev Methods Primers, 2021)
  \item Marees et al.\ -- A tutorial on conducting GWAS (step-by-step, very practical)
  \item LDSC wiki on GitHub --- the definitive usage guide
\end{itemize}
